% ===========================================================================
\section{Sub-Gaussian and sub-exponential}\label{sec:subgaussian}
\warmth{0}\quad\whisper{Turn ``what does the tail look like'' into a norm. Once classified, tail bounds come for free.}

\begin{strand}
In one line: rather than recompute MGFs each time, classify random variables by the
\emph{shape of their tail}. Sub-Gaussian = tail no heavier than a Gaussian's; sub-exponential
= no heavier than an exponential's. Each family gets an Orlicz norm, and concentration
becomes arithmetic on those norms.
\end{strand}

\block{The paradigm}
\trigger{You want to handle a whole class of tails at once, and have the tail of sums and products follow \emph{mechanically}.}
Four equivalent characterizations (sub-Gaussian, write $\|X\|_{\psi_2}$): tail
$\Pr[|X|>t]\le 2e^{-t^2/K^2}$ $\Leftrightarrow$ moments $\|X\|_p\lesssim K\sqrt p$ $\Leftrightarrow$
MGF $\E e^{\lambda X}\le e^{K^2\lambda^2/2}$ (after centering) $\Leftrightarrow$ $\E e^{X^2/K^2}\le 2$.

\begin{strand}
Three-line heuristic: (1) \keyword{a centered bounded variable is sub-Gaussian} (Hoeffding's
lemma), so the entire ``bounded-difference'' family lands in range here; (2) sums of
sub-Gaussians stay sub-Gaussian and the proxy variances \emph{add}:
$\|\sum X_i\|_{\psi_2}^2\le\sum\|X_i\|_{\psi_2}^2$; (3) a product of two sub-Gaussians is sub-\emph{exponential},
which is exactly why covariance / quadratic-form estimates drop to the sub-exponential family.
\end{strand}

\block{Cheat-sheet: the two tail families}
\noindent\begin{tabularx}{\linewidth}{@{}l l L@{}}
\toprule
{\color{indigo}Theorem} & {\color{indigo}Setting} & \multicolumn{1}{l}{\color{indigo}Tail bound / when to use}\\
\midrule
Hoeffding lemma   & $\E X=0$, $X\in[a,b]$            & $\E e^{\lambda X}\le e^{\lambda^2(b-a)^2/8}$; bridge bounded $\to$ sub-Gaussian\\
Hoeffding ineq.\  & $X_i\in[a_i,b_i]$ independent    & $\Pr[|S-\E S|\ge t]\le 2e^{-2t^2/\sum(b_i-a_i)^2}$\\
General Hoeffding & $X_i$ sub-Gaussian               & $\Pr[|\sum a_iX_i|\ge t]\le 2e^{-ct^2/\|a\|_2^2K^2}$\\
Bernstein         & $X_i$ sub-exp., $\E X_i=0$       & $\Pr[|S|\ge t]\le 2\exp\!\big[-c\min(\tfrac{t^2}{\sum\nu_i^2},\tfrac{t}{\max b_i})\big]$\\
Bennett           & $|X_i|\le b$, variance $\sigma^2$& variance-aware; Gaussian for small $t$, Poisson for large $t$\\
\bottomrule
\end{tabularx}

\block{Main results \& proof skeletons}

\begin{definition}[Sub-Gaussian norm]\label{def:subg}
$\|X\|_{\psi_2}=\inf\{K>0:\ \E\exp(X^2/K^2)\le 2\}$. Call $X$ \keyword{sub-Gaussian} if this is finite.
(Sub-exponential $\|X\|_{\psi_1}$ replaces $X^2$ by $|X|$.)
\end{definition}

\begin{proposition}[Four equivalences]\label{prop:subg-equiv}
For centered $X$ the four characterizations above are pairwise equivalent, with constants differing only by absolute factors.
\end{proposition}
\begin{proof}
Skeleton: tail $\Rightarrow$ moments (integrate the tail) $\Rightarrow$ MGF (sum the series $\sum K^p p^{p/2}\lambda^p/p!$)
$\Rightarrow$ tail (Chernoff, §\ref{sec:chernoff}) $\Rightarrow$ Orlicz. \TODO{track the constants through the four steps; draw it as a cycle}
\end{proof}

\begin{lemma}[Hoeffding's lemma]\label{lem:hoeffding}
$\E X=0$, $X\in[a,b]$ a.s.\ $\Rightarrow \E e^{\lambda X}\le e^{\lambda^2(b-a)^2/8}$.
\end{lemma}
\begin{proof}
Skeleton: set $\psi(\lambda)=\log\E e^{\lambda X}$; compute $\psi(0)=\psi'(0)=0$ and
$\psi''(\lambda)=\Var_{\lambda}(X)\le(b-a)^2/4$ (the variance under the tilted measure is at most a quarter of the range squared); Taylor remainder.
\TODO{prove $\psi''\le(b-a)^2/4$ via the variance-range estimate}
\end{proof}

\begin{theorem}[Hoeffding's inequality]\label{thm:hoeffding}
$X_i\in[a_i,b_i]$ independent, $S=\sum X_i$. Then $\Pr[|S-\E S|\ge t]\le 2\exp\!\big(-\tfrac{2t^2}{\sum_i(b_i-a_i)^2}\big)$.
\end{theorem}
\begin{proof}\TODO{feed Lem~\ref{lem:hoeffding} into the master bound of §\ref{sec:chernoff} + Lem~\ref{lem:mgf-add}, optimize $\lambda$}\end{proof}

\begin{theorem}[Bernstein's inequality]\label{thm:bernstein}
$X_i$ independent, centered, sub-exponential. Then the tail is Gaussian for small $t$ and exponential for large $t$ (see the cheat-sheet row).
\end{theorem}
\begin{proof}
Skeleton: the single-term sub-exponential MGF bound $\E e^{\lambda X_i}\le e^{\nu_i^2\lambda^2/2}$ holds only for
$|\lambda|<1/b_i$, so the choice of $\lambda$ is constrained by $1/b$, producing the two-regime $\min$.
\TODO{do the constrained $\lambda$-optimization (the boundary between the two regimes)}
\end{proof}

\begin{strand}
Why it goes this way: a Gaussian tail comes from ``MGF dominated everywhere by $e^{c\lambda^2}$'';
an exponential tail comes from ``dominated only for $|\lambda|<1/b$''. The constrained $\lambda$-optimization
automatically yields Bernstein's $\min(t^2/\nu,\,t/b)$ split, which is the whole secret of being \emph{variance-aware}.
\end{strand}

\begin{example}[Johnson–Lindenstrauss, preview]\label{ex:jl}
The squared length of a Gaussian random projection $\frac1{\sqrt m}Gx$ is a $\chi^2$-type sum (sub-exponential),
and Bernstein gives $\Pr[|\,\|\Pi x\|^2-\|x\|^2\,|>\epsilon\|x\|^2]\le 2e^{-cm\epsilon^2}$.
\TODO{add the union bound giving $m=O(\epsilon^{-2}\log N)$}
\end{example}

\loose{Check the constant: does Hoeffding's $2t^2/\sum(b-a)^2$ agree with Appendix A's A.1.18 (range $\le1$ gives $e^{-2a^2/n}$)?}%
\pick{mgf}\quad
\warp{subgaussian} The sub-Gaussian thread is reused throughout §\ref{sec:bdd} (McDiarmid) and §\ref{sec:matrix} (matrix Bernstein) via \pick{subgaussian}.
